\chapter{Results}
\label{chap:results}

This chapter presents the empirical findings derived from the simulation runs outlined in the experimental design (\ref{chap:experimental_design}). By analyzing the telemetry data extracted from the GeoMAS engine, this section validates the framework's capability to model complex geopolitical dynamics and addresses the four primary Research Questions. The analysis progresses from evaluating the macro-stability of the baseline system to dissecting the micro-cognitive behaviors of individual agents when subjected to exogenous shocks or counterfactual constraints.

% ============================================================
% WORKFLOW COMPLESSIVO — LEGGI PRIMA DI TUTTO
% ============================================================
% STEP 0 — Selezione baseline
%   Hai eseguito 3 run parallele identiche (stesso map_seed, history_seed, n_nations=8).
%   Visualizza le 3 simulazioni nel notebook (Section 1.5 Global KPIs Summary).
%   Confronta le metriche aggregate (avg deception, avg satisfaction, war rate) tra i 3.
%   Scegli come BASELINE_ID quella con comportamento più rappresentativo della mediana
%   (non la più outlier). Annota il simulation_id scelto.
%
% STEP 1 — Esegui i fork degli scenari dalla baseline scelta:
%   Dalla dashboard (Analysis Mode), carica il turno ~T=25 della baseline.
%   Esegui 3 fork con scenario diverso ciascuno (Pandemic, Insurrection, Resource Discovery).
%   Annota i simulation_id dei fork.
%
% STEP 2 — Esegui i fork counterfactuali XAI:
%   Dalla dashboard, carica il turno di maggior tensione della baseline (es. T dove
%   la nazione più aggressiva ha guerra attiva). Inietta un constraint XAI
%   (es. FORCE ATTACK su una nazione pacifica, oppure FORBID NUCLEAR sulla più aggressiva).
%   Annota il simulation_id del fork counterfactuale.
%
% STEP 3 — Analizza nel notebook selezionando baseline_id e i fork_id per ciascun capitolo.
%   Esporta i plot come .png e inseriscili in results/figures/.
%
% STEP 4 — Compila le sezioni di questo file nell'ordine in cui appaiono.
% ============================================================

\section{Baseline Telemetry and System Stability}
\label{sec:res_baseline}
\textit{Addresses RQ4: Emergent Socio-Political Equilibria.}

This section analyzes the unperturbed baseline simulations to establish the system's default behavioral attractors, ensuring that the engine produces coherent macro-dynamics without spiraling into premature chaos or stagnation.

\paragraph{Agent Cognitive Quality.}
Figure~\ref{fig:baseline_deception_coherence} reports the global average Deception Score and Coherence Score across all 50 turns of the baseline simulation. Both metrics exhibit striking temporal stability: the Coherence Score stabilizes immediately at approximately $0.87$ and remains in the range $[0.82, 0.95]$ throughout the entire run, while the Deception Score converges to approximately $0.27$ within the first two turns and shows no discernible trend or crisis-induced spikes. The aggregate summary (Table~\ref{tab:baseline_kpis}) confirms these observations, reporting a mean Coherence of $0.85$ and a mean Deception of $0.27$ over the full simulation horizon.

The absence of an accelerating deception trajectory is theoretically significant: in a system driven to conflict by resource scarcity, one might expect agents to progressively decouple their public rhetoric from their private strategic intent. The observed stability instead suggests that LLM agents maintain a consistent strategic personality throughout the run, with deception bounded by ideological constraints embedded in their respective system prompts. The finer breakdown of deception by Government Type and Global Strategy is examined in Section~\ref{sec:res_deception}.

\begin{figure}[!htb]
    \centering
    \includegraphics[width=\linewidth]{figure/baseline_deception_coherence.png}
    \caption{Global average Deception Score (red) and Coherence Score (green) over 50 turns. Both metrics are stable throughout the simulation horizon, with no crisis-induced trend or spike.}
    \label{fig:baseline_deception_coherence}
\end{figure}

\paragraph{Resource Allocation Dynamics: A Three-Phase Geopolitical Cycle.}
The Guns vs.\ Butter chart (Figure~\ref{fig:baseline_guns_butter}) reveals a more complex dynamic than simple equilibrium by exposing three structurally distinct phases in the joint evolution of global Trade Volume and Units Created.

During the \textit{initialization phase} ($T = 1$--$7$), nations deploy initial military stocks from pre-allocated resources, producing a large pulse of unit creation ($\approx\!200{,}000$ units in turn 1) which declines to a steady-state production rate of approximately $32{,}000$--$42{,}000$ units per turn by turn 7. Trade Volume grows sharply from zero to approximately $400{,}000$, reflecting the formation of early diplomatic and commercial agreements.

The \textit{competitive equilibrium phase} ($T = 8$--$18$) is characterized by oscillating trade volumes in the range $[220{,}000, 410{,}000]$ and stable military production. This alternating pattern reflects turn-to-turn variation as nations balance trade against resource hoarding in response to changing diplomatic tensions --- a behavioral signature consistent with the mixed population of \textsc{Coalition Builder} and \textsc{Total Expansionism} strategies in the baseline.

The most notable emergent phenomenon is the \textit{economic paralysis phase} ($T = 19$--$48$): both Trade Volume and Units Created collapse to near-zero and remain depressed for approximately 30 turns. Cross-referencing with the relationship matrix stored in the simulation snapshots confirms that this phase coincides with AGRIA's simultaneous declaration of war against multiple adversaries. During active multi-front warfare, trade embargoes eliminate inter-nation commerce while war expenditures deplete national budgets to the point where new unit production becomes financially infeasible. This constitutes an emergent \emph{war-driven economic paralysis} --- the system self-generates a Nash equilibrium of mutual attrition with no economically productive activity. The recovery spike at $T = 49$--$50$ corresponds to the resolution of hostilities and the resumption of pent-up trade, confirming that the paralysis is war-contingent rather than a persistent structural condition.

\begin{figure}[!htb]
    \centering
    \includegraphics[width=\linewidth]{figure/baseline_guns_butter.png}
    \caption{Guns vs.\ Butter chart for the baseline simulation. Trade Volume (blue, left axis) and Units Created (orange, right axis) reveal three distinct geopolitical phases: initialization ($T=1$--$7$), competitive equilibrium ($T=8$--$18$), and war-driven economic paralysis ($T=19$--$48$), followed by a recovery spike at $T=49$--$50$.}
    \label{fig:baseline_guns_butter}
\end{figure}

\paragraph{Power Distribution and Hegemonic Structure.}
Figure~\ref{fig:baseline_power_projection} illustrates the Power Projection trajectory of all eight nations across the full simulation. The distribution is markedly hierarchical and asymmetric. AGRIA emerges as the uncontested hegemon, accumulating power monotonically from approximately $15{,}000$ at $T=1$ to $47{,}000$ by $T=50$ --- a trajectory consistent with an aggressive \textsc{Total Expansionism} strategic profile that prioritizes territorial acquisition. A secondary tier comprising OSTER and MERIDIA at approximately $35{,}000$--$38{,}000$ achieves stable growth, suggesting sustained resource accumulation and effective mutual-defense alliances. KRELL maintains a flat profile at approximately $24{,}000$, consistent with an \textsc{Armed Isolationism} posture that preserves existing strength without active expansion. NORD and SYLVAN exhibit moderate growth to $14{,}000$ and $13{,}000$ respectively, while VALKYR remains marginal at approximately $9{,}000$.

The most notable trajectory is ZENTORA, whose Power Projection collapses from approximately $10{,}000$ at $T=1$ to near zero by $T=10$, remaining permanently diminished for the remainder of the simulation. This trajectory indicates a catastrophic early defeat --- consistent with conquest by AGRIA in the initialization phase --- from which ZENTORA never recovers territorial or economic strength. This outcome directly addresses RQ4: the baseline system does \emph{not} converge toward a multipolar balance of power; instead, it self-organizes around a dominant hegemon, a partially co-opted secondary tier, and peripheral marginal states --- a configuration consistent with classical power transition theory.

\begin{figure}[!htb]
    \centering
    \includegraphics[width=\linewidth]{figure/baseline_power_projection.png}
    \caption{Power Projection evolution per nation over 50 turns. AGRIA emerges as the unchallenged hegemon ($\approx\!47{,}000$ by $T=50$), while ZENTORA collapses to near zero by $T=10$. The remaining nations stratify into stable mid-level and peripheral tiers.}
    \label{fig:baseline_power_projection}
\end{figure}

\paragraph{Global KPI Summary.}
Table~\ref{tab:baseline_kpis} reports the aggregate performance indicators for the full baseline run. Several observations merit attention. First, the Violence Rate of $0.01$ --- defined as the ratio of units destroyed to units created --- indicates that warfare produces negligible annihilation relative to military buildup: the 108 recorded territory changes occur primarily through overwhelming force superiority rather than symmetric attrition. Second, the average Public Satisfaction of $73.23$ confirms that the system maintains social stability even during prolonged conflict. Third, the complete absence of Civil Unrest events ($0.00$) across all 50 turns is particularly significant: despite resource pressures and multi-front wars, no nation crosses the unrest threshold, validating the robustness of the welfare mechanics in the GeoMAS engine.

\begin{table}[!htb]
    \centering
    \caption{Baseline Simulation --- Global KPI Summary (50 turns, 8 nations).}
    \label{tab:baseline_kpis}
    \begin{tabular}{lr}
        \hline
        \textbf{Metric} & \textbf{Value} \\
        \hline
        Total Territory Changes         & 108 \\
        Total Units Created             & $1{,}735{,}800$ \\
        Total Units Destroyed           & $22{,}968$ \\
        Total Trade Volume              & $17{,}286{,}392$ \\
        Avg.\ Deception Score           & 0.27 \\
        Avg.\ Coherence Score           & 0.85 \\
        Avg.\ Public Satisfaction       & 73.23 \\
        Violence Rate (destr./created)  & 0.01 \\
        Total Civil Unrest Turn-Nations & 0.00 \\
        \hline
    \end{tabular}
\end{table}
% NOTE TO AUTHOR: the spike in Units Created at T≈45 visible in the graph should be
% investigated by loading those specific turns in the dashboard (Envelopes tab) to
% determine which nation triggered a mass-mobilization and why (post-war recovery?
% arms race response to AGRIA hegemony?).

\paragraph{Diplomatic Network Evolution and Alliance Formation.}
Figure~\ref{fig:baseline_network} traces the evolution of the diplomatic relationship network at five equidistant snapshot turns. At $T=1$, the network is fully connected with all edges in \textsc{Peace} (gray), reflecting the initial tabula rasa in which no diplomatic history has been established. By $T=12$, the first \textsc{War} (red) and \textsc{Mutual Defense} (green) relationships appear, signaling the onset of active geopolitical competition. Over subsequent snapshots, the graph progressively partitions: green alliance edges cluster among the mid-tier powers while red war edges radiate outward from AGRIA toward its adversaries.

By $T=50$, the algorithm detects five structurally disconnected blocs:

\begin{itemize}
    \item \textbf{Bloc 1 — Anti-Hegemonic Coalition:} \{KRELL, MERIDIA, OSTER, SYLVAN\}. Four nations bound by \textsc{Mutual Defense} pacts form a coordinated defensive alliance against AGRIA's expansion. This is the dominant structural outcome and directly corresponds to classical balance-of-power theory: intermediate powers aggregate in response to a rising hegemon.
    \item \textbf{Bloc 2 — Isolated Hegemon:} \{AGRIA\}. Despite (or because of) its military dominance, AGRIA ends the simulation with no formal allies. Its \textsc{Total Expansionism} strategy, which optimizes for territorial acquisition, renders sustained alliance-building incoherent with its behavioral profile. The largest military power is diplomatically isolated.
    \item \textbf{Blocs 3--5 — Peripheral Isolates:} \{ZENTORA\}, \{VALKYR\}, \{NORD\}, each unable to attract alliance partners, whether because of resource exhaustion (ZENTORA), marginal power projection (VALKYR), or strategic disconnect (NORD).
\end{itemize}

This five-bloc configuration represents a unipolar system with organized peripheral resistance rather than a stable multipolar equilibrium, confirming RQ4's central hypothesis about the system's attractor states.

\begin{figure}[!htb]
    \centering
    \includegraphics[width=\linewidth]{figure/baseline_network_evolution.png}
    \caption{Diplomatic network evolution at five snapshot turns ($T=1, 12, 25, 37, 50$). Colors encode relationship state: green = \textsc{Mutual Defense}, red = \textsc{War}, yellow = \textsc{Non-Aggression}, gray = \textsc{Peace}. The network transitions from a fully peaceful initial state to a structured bipolar configuration with an isolated hegemon (AGRIA) opposed by a four-nation defensive coalition.}
    \label{fig:baseline_network}
\end{figure}

\paragraph{Socio-Political Stability and Satisfaction Divergence.}
Figure~\ref{fig:baseline_satisfaction} reports the public satisfaction trajectory of all eight nations alongside the global satisfaction dispersion ($\sigma$) and the civil unrest count per turn. Three observations stand out.

First, AGRIA's satisfaction grows nearly monotonically from ${\approx}75$ at $T=1$ to ${\approx}98$ by $T=50$. This is the satisfaction dividend of hegemonic victory: conquered territories provide resources and security that consistently exceed the welfare needs of the expanding population. Second, the majority of nations exhibit a characteristic sawtooth oscillation in the range $[55, 70]$, reflecting the alternating cycle of resource investment and satisfaction decay. No nation breaches the civil unrest threshold (dashed red line at 10) at any turn, and the Civil Unrest count remains identically zero throughout the entire simulation — confirming the aggregate value reported in Table~\ref{tab:baseline_kpis}.

Third, and arguably most theoretically significant, the satisfaction dispersion $\sigma$ grows from approximately $8$ at $T=1$ to approximately $22$ by $T=50$, representing a near-threefold increase in inter-nation inequality. This widening gap between AGRIA's near-perfect satisfaction and the stagnating peripheral nations constitutes an emergent \emph{winner-takes-all divergence}: the system self-generates increasing socio-economic inequality even in the absence of any externally imposed redistribution mechanism.

\begin{figure}[!htb]
    \centering
    \includegraphics[width=\linewidth]{figure/baseline_satisfaction.png}
    \caption{Left: public satisfaction per nation over 50 turns. AGRIA grows monotonically toward 100; most nations oscillate in $[55, 70]$; ZENTORA collapses toward the strike threshold by $T=10$. Center: satisfaction dispersion ($\sigma$) grows from 8 to 22, indicating increasing inter-nation inequality. Right: civil unrest count per turn — identically zero throughout the simulation.}
    \label{fig:baseline_satisfaction}
\end{figure}

\paragraph{Territorial Dynamics.}
Figure~\ref{fig:baseline_territory} reports the cumulative count of provinces that change ownership across the 50-turn baseline. The trajectory is approximately linear, accumulating 108 territorial transfers at a mean rate of ${\approx}2.2$ provinces per turn. Notably, the slope is slightly steeper in the early phase ($T=1$--$20$), reflecting AGRIA's rapid conquest of ZENTORA and initial territorial expansion, then moderates as the remaining nations consolidate defensive positions. The linearity of the curve indicates that geopolitical expansion is a sustained, steady-state process rather than a punctuated equilibrium of discrete crises — a behavioral signature consistent with a hegemon engaged in continuous, resource-driven territorial absorption rather than apocalyptic blitzkrieg.

\begin{figure}[!htb]
    \centering
    \includegraphics[width=\linewidth]{figure/baseline_territory_changes.png}
    \caption{Cumulative territory changes (provinces changing ownership) over 50 turns. The near-linear trajectory accumulates 108 total transfers at a mean rate of ${\approx}2.2$ per turn, consistent with AGRIA's sustained territorial expansion strategy.}
    \label{fig:baseline_territory}
\end{figure}

\paragraph{Engine Governance: Action Acceptance and Presidential Oversight.}
Tables~\ref{tab:engine_acceptance} and~\ref{tab:veto_rate} report the engine-level action acceptance rates and presidential veto rates respectively, providing a quantitative health check of the GeoMAS decision pipeline.

The overall engine acceptance rate is high, with five of eight nations achieving rates at or above $97.6\%$. AGRIA's single rejected action ($99.6\%$) confirms that its aggressive strategy is computationally well-specified and respected by the action validation layer. The two notable exceptions are ZENTORA ($89.6\%$) and KRELL ($76.6\%$). ZENTORA's elevated failure rate is directly attributable to catastrophic resource exhaustion: a nation reduced to near-zero provinces and budget attempts actions its economy cannot support, such as large-scale military deployments on empty reserves. KRELL's $76.6\%$ rate, by contrast, likely reflects the \textsc{Armed Isolationism} strategy's tendency to propose self-sustaining defensive arrangements that the engine rejects as diplomatically inconsistent with its relationship state. Both cases are behaviorally interpretable rather than indicative of system malfunction.

The presidential veto rate reveals a complementary dimension of the decision hierarchy. ZENTORA exhibits by far the highest veto rate ($28.3\%$), consistent with a head of state in crisis governance mode who rejects ministerial proposals perceived as strategically untenable given the nation's diminished capacity. AGRIA's low veto rate ($5.0\%$) confirms strong presidential alignment with the expansionist military agenda. Most notably, MERIDIA and SYLVAN achieve $0.0\%$ veto rates — complete presidential alignment with ministerial proposals — consistent with their stable mid-tier positions and coherent \textsc{Coalition Builder} and \textsc{Democracy} profiles where domestic consensus is the dominant behavioral attractor.

\begin{table}[!htb]
    \centering
    \caption{Per-Nation Engine Action Acceptance Rate (baseline, full 50-turn run).}
    \label{tab:engine_acceptance}
    \begin{tabular}{lrrrr}
        \hline
        \textbf{Nation} & \textbf{Failed} & \textbf{Success} & \textbf{Total} & \textbf{Rate (\%)} \\
        \hline
        MERIDIA & 0  & 248 & 248 & 100.0 \\
        OSTER   & 0  & 227 & 227 & 100.0 \\
        VALKYR  & 0  & 200 & 200 & 100.0 \\
        AGRIA   & 1  & 232 & 233 & 99.6  \\
        NORD    & 3  & 158 & 161 & 98.1  \\
        SYLVAN  & 6  & 243 & 249 & 97.6  \\
        ZENTORA & 7  & 60  & 67  & 89.6  \\
        KRELL   & 30 & 98  & 128 & 76.6  \\
        \hline
    \end{tabular}
\end{table}

\begin{table}[!htb]
    \centering
    \caption{Per-Nation Presidential Veto Rate (baseline, full 50-turn run).}
    \label{tab:veto_rate}
    \begin{tabular}{lrrrr}
        \hline
        \textbf{Nation} & \textbf{APPROVE} & \textbf{VETO} & \textbf{Total} & \textbf{Veto Rate (\%)} \\
        \hline
        ZENTORA & 33  & 13 & 46  & 28.3 \\
        OSTER   & 127 & 16 & 143 & 11.2 \\
        VALKYR  & 100 & 7  & 107 & 6.5  \\
        AGRIA   & 133 & 7  & 140 & 5.0  \\
        NORD    & 61  & 2  & 63  & 3.2  \\
        KRELL   & 146 & 2  & 148 & 1.4  \\
        MERIDIA & 148 & 0  & 148 & 0.0  \\
        SYLVAN  & 149 & 0  & 149 & 0.0  \\
        \hline
    \end{tabular}
\end{table}

\section{Macrodynamics: Scenario Impact Analysis}
\label{sec:res_scenario}
\textit{Addresses RQ1: Impact of Exogenous Scenario Perturbations.}

This section evaluates the resilience and adaptability of the multi-agent system when subjected to global environmental shocks, including insurrection, pandemic, and resource discovery events.

% ── NOTEBOOK: usa BASELINE_ID vs uno dei fork-scenario alla volta. ───────────
%
% [§2.1] Multi-Metric Overlay (Baseline vs Scenario)
%   → Seleziona BASELINE_ID come sim A e il fork con scenario come sim B.
%   → Metriche chiave da sovrapporre: public_satisfaction, trade_volume,
%     soldiers (unità militari), budget medio, global_deception_avg.
%   → SCRIVI per OGNI scenario (pandemia, insurrezione, resource discovery):
%     - "A quale turno l'effetto dello shock diventa visibile nelle metriche?"
%     - "Quali nazioni reagiscono per prime? Come?"
%     - "Il sistema si riprende (resilienza) o diverge permanentemente dalla baseline?"
%
% [§2.2] Trust Heatmap Comparison
%   → Confronta la trust matrix BASELINE vs SCENARIO a T_fork, T_fork+10, T_fine.
%   → SCRIVI: lo scenario altera le relazioni diplomatiche? La pandemia, per esempio,
%     porta a riduzione del trust globale (competizione per risorse) o a solidarietà?
%     L'insurrezione genera un cambio di alleanze verso la nazione colpita?
%
% [§2.3] Summary Statistics (Δ Baseline → Scenario)
%   → Tabella con i delta: Δ satisfaction, Δ trade_volume, Δ wars, Δ territories.
%   → INSERISCI questa tabella come "Table X: Scenario Impact Summary".
%   → SCRIVI: quale scenario ha l'impatto quantitativamente maggiore?
%     Quale il sistema assorbe più facilmente?
%
% ANALISI QUALITATIVA (da fare nella dashboard Streamlit, non nel notebook):
%   → Carica il turno T_fork+1 di ogni simulazione fork nella dashboard.
%   → Leggi i "Notable Events" e gli "Envelopes" delle nazioni più colpite.
%   → CITA 1-2 esempi di ragionamento testuale dell'agente che risponde allo shock.
%     Es: "In risposta all'insurrection, AGRIA nel turno T+2 scrive nel suo
%     Defense Reasoning: '...'. Questo dimostra che l'agente percepisce la
%     destabilizzazione e adatta la strategia militare."
%
% TESTO DA SCRIVERE (struttura paragrafi):
%   Un sottoparagrafo per ogni scenario testato:
%   \subsection*{Scenario: [Nome]}
%   P1: Descrivi lo shock e il turno di iniezione.
%   P2: Analisi quantitativa (Δ metriche, overlay plot).
%   P3: Analisi qualitativa (come gli agenti ragionano, citazione dal log).
%   P4: Conclusione: il sistema è resiliente o diverge? Risposta calibrata?

\section{Microdynamics: Strategic Deception and Moral Washing}
\label{sec:res_deception}
\textit{Addresses RQ3: Alignment of Diplomatic Intent and Action.}

This section investigates the internal cognitive processes of the LLM agents, specifically assessing their propensity to decouple public rhetoric from private strategic intent based on their Government Type and Global Strategy.

% ── NOTEBOOK: usa BASELINE_ID per tutta questa sezione. ─────────────────────
%
% [§4.1] Deception by Domain (Defense vs Foreign)
%   → Il deception_defense misura se le azioni militari private divergono dalla
%     dichiarazione pubblica. Il deception_foreign misura lo stesso per la diplomazia.
%   → SCRIVI: quale dominio mostra più inganno? È coerente con l'aspettativa teorica
%     (le decisioni militari dovrebbero essere più occultate di quelle diplomatiche)?
%
% [§4.2] Deception Breakdown by GovernmentType
%   → Boxplot: AUTHORITARIAN vs DEMOCRACY vs THEOCRACY.
%   → SCRIVI: gli agenti AUTHORITARIAN hanno deception score sistematicamente più alta?
%     Questo conferma che il contesto ideologico nel prompt influenza il behavior
%     del modello LLM, producendo agenti con personalità distinte.
%
% [§4.3] Deception Breakdown by GlobalStrategy
%   → Boxplot: TOTAL_EXPANSIONISM vs COALITION_BUILDER vs SCORCHED_EARTH vs ARMED_ISOLATIONISM.
%   → SCRIVI: TOTAL_EXPANSIONISM dovrebbe mostrare più deception (pianificano aggressione
%     ma la dichiarano come "difesa del confine"). COALITION_BUILDER dovrebbe avere
%     deception più bassa (interesse nel mantenere coerenza per costruire alleanze).
%
% [§4.4] Deception × Coherence Scatter
%   → Osserva se c'è correlazione negativa (più inganno = meno coerenza).
%   → SCRIVI: la correlazione è attesa teoricamente (un agente che mente pubblicamente
%     deve anche razionalizzare internamente). La sua presenza empirica valida il
%     funzionamento della pipeline di scoring.
%
% [§4.5] Public Statement Text Mining
%   → Frequenza keyword "liberation", "sacred", "peace", "necessary", "defense",
%     "sovereignty", "democracy", "stability" etc. nei public statement.
%   → SCRIVI: le nazioni AUTHORITARIAN usano significativamente più keyword di
%     moral washing (liberation, defense, sacred) rispetto alle DEMOCRACY?
%     Questo è l'evidenza empirica del fenomeno "Moral Washing".
%
% [§4.6] Public vs Private Intent Examples (High-Deception Turns)
%   → Filtra i 5 turni con deception score più alta nella baseline.
%   → Per ciascuno, mostra il public_statement e il defense_private_reasoning fianco a fianco.
%   → INSERISCI come "Exhibit X: High-Deception Turn Example" con una box formattata.
%   → SCRIVI: "In Turn T, Nation X declared publicly: '...' [high moral framing].
%     Simultaneously, its private military reasoning read: '...' [explicit resource grab].
%     This constitutes a canonical instance of Moral Washing."
%
% TESTO DA SCRIVERE (struttura paragrafi):
%   P1: Panoramica del Deception Score medio nella baseline (valore numerico da §1.1).
%   P2: Analisi per Government Type — gli AUTHORITARIAN ingannano di più? (§4.2)
%   P3: Analisi per Strategy — TOTAL_EXPANSIONISM vs COALITION_BUILDER (§4.3).
%   P4: Evidenza testuale del Moral Washing (§4.5 + §4.6 examples).
%   P5: Conclusione: il LLM spontaneamente disaccoppia retorica pubblica da intent privato.

\section{Causal Impacts: Counterfactual Analysis and XAI}
\label{sec:res_counterfactual}
\textit{Addresses RQ2: Causal Impact of Forced Strategic Decisions.}

Utilizing the state-forking mechanism, this section traces the causal chain of localized cognitive injections, demonstrating how a single behavioral constraint imposed on one agent cascades through the global geopolitical equilibrium.

% ── NOTEBOOK: usa BASELINE_ID (sim A) vs il fork XAI (sim B). ───────────────
%
% PREPARAZIONE — scegli il fork point e il constraint:
%   → Identifica nella baseline il turno più "critico" (es. il turno in cui
%     una nazione democratica ha le risorse più basse e ha firmato una pace recente).
%   → Scegli come injection: FORCE ATTACK (da nazione X verso Y) oppure
%     FORBID NUCLEAR (sulla nazione più armata).
%   → Annota: fork_turn T, nation_id iniettata, constraint type, nation target.
%
% [§3.2] Prompt Audit: Envelope Comparison at Fork Point
%   → Seleziona BASELINE_ID e XAI_FORK_ID, turno T (fork turn).
%   → Il notebook mostra side-by-side il defense_private_reasoning della nazione
%     iniettata nella baseline e nel fork.
%   → ESPORTA questo confronto come immagine / tabella formattata.
%   → SCRIVI: "At Turn T, the injected constraint replaced the agent's natural
%     reasoning with: '[FORCED: ATTACK ...]'. The agent's observable output JSON
%     — the action selected — changed from [no_attack] to [MILITARY_ATTACK].
%     This demonstrates the mechanical traceability of the XAI injection."
%
% [§3.1] Divergence Visualization
%   → Grafico overlay Baseline vs XAI Fork su: satisfaction, trust medio, unità militari,
%     territori conquistati, trade volume.
%   → SCRIVI: "From Turn T onward, the trajectories visibly diverge. By Turn T+5,
%     [Metric X] in the XAI Fork is [N]% higher/lower than the baseline. By T+15,
%     [Metric Y] shows [evidence of cascade effect]."
%   → Identifica l'effetto a cascata: la vittima dell'ATTACK forzato chiede alleanze?
%     Le altre nazioni cambiano trust verso l'aggressore?
%
% [§3.3] Cumulative Divergence Table
%   → Tabella con la somma cumulativa delle divergenze per metrica.
%   → INSERISCI come "Table X: Cumulative Divergence from Counterfactual Injection".
%   → SCRIVI: quale metrica accumula la divergenza più grande nel tempo?
%     Questo risponde a RQ2: quale dimensione del sistema è più sensibile all'injection?
%
% ANALISI QUALITATIVA (da fare nella dashboard):
%   → Carica i turni T+1, T+5, T+10 del fork XAI nella dashboard.
%   → Osserva nelle Nations View come la nazione bersaglio dell'attacco forzato
%     cambia la sua strategia nelle request successive.
%   → CITA un esempio: "In the subsequent turn, Nation Y — having been attacked — 
%     shifted its Foreign Minister's reasoning from cooperative to defensive:
%     '...' [cita reasoning]. Its trust toward Nation X dropped from [X] to [Y]."
%
% TESTO DA SCRIVERE (struttura paragrafi):
%   P1: Descrivi il fork point, la nazione iniettata, il constraint scelto e perché.
%   P2: Prompt Audit — mostra il cambio di reasoning all'injection turn (§3.2).
%   P3: Analisi quantitativa della divergenza nelle metriche globali (§3.1 + §3.3).
%   P4: Effetto a cascata — come reagiscono le altre nazioni a questa singola decisione?
%   P5: Conclusione: il forking + injection dimostra la causalità. La singola deviazione
%       è tracciabile attraverso tutto il sistema, validando il framework come strumento XAI.

\section{Technical Performance and Computational Overhead}
\label{sec:res_performance}
This brief section reports on the operational efficiency, token utilization, and estimated costs associated with running the DeepSeek V3.2 models asynchronously over extended simulation horizons.

% ── NOTEBOOK: usa BASELINE_ID. ───────────────────────────────────────────────
%
% NON C'È UN CAPITOLO DEDICATO AL TOKEN COST nel notebook attuale.
% Usa i dati dal database direttamente:
%   SELECT model, SUM(total_tokens), SUM(cost), COUNT(*) FROM token_usage
%   WHERE simulation_id = '<BASELINE_ID>'
%   GROUP BY model, agent_type
%
% Metriche da riportare:
%   - Totale token consumati per run da 50 turni.
%   - Costo stimato in USD (calcolato da litellm o dal campo cost nel DB).
%   - Breakdown per agent_type (defense, foreign, economy, opinion, president).
%   - Latenza media per turno (se loggata).
%   - Numero di chiamate LLM totali per run.
%
% SCRIVI:
%   P1: "A full 50-turn baseline simulation with 8 nations consumed approximately
%       [X] total tokens at an estimated cost of $[Y] USD using DeepSeek V3.2
%       via the LiteLLM gateway."
%   P2: "The Defense and Foreign minister agents are the primary token consumers,
%       accounting for [Z]% of total usage, consistent with their higher
%       reasoning complexity (multi-step JSON generation + private reasoning field)."
%   P3: "The asynchronous ministerial proposal pipeline (concurrent per-nation) reduced
%       wall-clock time per turn from ~[A]s (sequential) to ~[B]s (parallel),
%       enabling the full 50-turn run within [C] hours."
%
% TABELLA SUGGERITA:
%   \begin{table}[h]
%     \caption{Token Usage and Cost Breakdown — 50-turn Baseline}
%     Agent Type | Total Tokens | Avg per Turn | Cost (USD)
%     Defense    | ...          | ...          | ...
%     Foreign    | ...          | ...          | ...
%     Economy    | ...          | ...          | ...
%     Opinion    | ...          | ...          | ...
%     President  | ...          | ...          | ...
%     TOTAL      | ...          | ...          | ...
%   \end{table}